\documentclass{article}
\usepackage{amsmath,amssymb}
\begin{document}
% (1) 危险率与区间概率
\begin{equation}
\lambda_{e,k}=a_e[\max(v_{e,k}/v_e^0-1,0)]^{b_e},\quad p_{e,k}=1-e^{-\lambda_{e,k}\Delta t}
\tag{1}
\end{equation}
% (2) 累计首次故障概率
\begin{equation}
F_e(t)=1-\exp[-\sum_{k:t_k<t}\lambda_{e,k}\Delta t]
\tag{2}
\end{equation}
% (3) 预测加权恢复面积
\begin{equation}
R={A_{\rm past}+\sum_s\pi_s\sum_{k=k_0}^{N-1}\Delta t\sum_i(w_c\ell^c_{i,k,s}+w_g\ell^g_{i,k,s})\over A_{\rm demand}}
\tag{3}
\end{equation}
% (4) 实际失供电量
\begin{equation}
\mathrm{EENS}=\sum_{k=0}^{N-1}\Delta t\sum_i(d^c_{i,k}+d^g_{i,k}-\ell^c_{i,k}-\ell^g_{i,k})
\tag{4}
\end{equation}
% (5) SOC跨期动态
\begin{equation}
E_{j,k+1}=E_{j,k}+\Delta t(\eta^c_jP^{\rm ch}_{j,k}-P^{\rm dis}_{j,k}/\eta^d_j)
\tag{5}
\end{equation}
% (6) MGO接口净功率
\begin{equation}
p_{j,k}=P^{\rm PV}_{j,k}+P^{\rm DG}_{j,k}+P^{\rm dis}_{j,k}-P^{\rm ch}_{j,k}-\ell^c_{j,k}-\ell^g_{j,k}
\tag{6}
\end{equation}
% (7) 条件非预期性
\begin{equation}
\mathcal I_{s,k}=\mathcal I_{r,k}\ \Longrightarrow\ u_{s,k}=u_{r,k}
\tag{7}
\end{equation}
% (8) 实体任务连续性
\begin{equation}
q_k\ne\varnothing,\ r_k>0\ \Longrightarrow\ q_{k+1}=q_k\quad\text{until completion}
\tag{8}
\end{equation}
% (9) 有限规则类的最优值界
\begin{equation}
L_b=\max_{p\in\mathcal P}L_{b,p},\quad U_b=\max_{p\in\mathcal P}U_{b,p}
\tag{9}
\end{equation}
% (10) 启动等待差值与区间
\begin{equation}
D_R=R^*_{\rm start}-R^*_{\rm wait}\in[L_{\rm start}-U_{\rm wait},\ U_{\rm start}-L_{\rm wait}]
\tag{10}
\end{equation}
% (11) 同模型追加求解界合并
\begin{equation}
L^{m+1}=\max(L^m,\widehat L^{m+1}),\quad U^{m+1}=\min(U^m,\widehat U^{m+1})
\tag{11}
\end{equation}
% (12) 后向管终端集合
\begin{equation}
\mathcal B_N=[E_{\rm terminal},E_{\max}]
\tag{12}
\end{equation}
% (13) 区间前驱递推
\begin{equation}
\mathcal B_k=[E_{\min},E_{\max}]\cap[\underline B_{k+1}-\delta_{\max},\overline B_{k+1}-\delta_{\min}]
\tag{13}
\end{equation}
% (14) 观测前先交可能子分支
\begin{equation}
\mathcal J_{g,k+1}=\bigcap_{h\in\operatorname{child}(g)}\mathcal B_{h,k+1},\quad\mathcal B_{g,k}=\mathcal X\cap\operatorname{Pre}(\mathcal J_{g,k+1})
\tag{14}
\end{equation}
% (15) 共同系数的全轨迹内近似
\begin{equation}
v=\sum_r\lambda_rv^r,\quad x=\sum_r\lambda_rx^r,\quad\lambda_r\ge0,\ \sum_r\lambda_r=1
\tag{15}
\end{equation}
% (16) 主配反送路径贡献的独立对照
\begin{equation}
\Delta R_{\rm path}=R_{\rm allow\ upward}^{\rm actual}-R_{\rm no\ upward}^{\rm actual}
\tag{16}
\end{equation}
% (17) 含有限盒余项的拉格朗日界
\begin{equation}
g(\lambda,\mu)=b^T\lambda+d^T\mu+\sum_j\min(r_jl_j,r_ju_j),\quad r=c-A^T\lambda-C^T\mu,\ \lambda\le0
\tag{17}
\end{equation}
% (18) 完整轨迹列进入限制主问题
\begin{equation}
z=V\lambda,\quad A_yy+A_zV\lambda\le b,\quad l_z\le V\lambda\le u_z,\quad \mathbf{1}^T\lambda=1,\ \lambda\ge0
\tag{18}
\end{equation}
% (19) 保持终端条件的前缀可延续性
\begin{equation}
n_{\rm fail}=\min\{n:\nexists x\in\mathcal F,\ p_{:n}=\hat p_{:n},\ q_{:n}=\hat q_{:n}\}
\tag{19}
\end{equation}
% (20) 最小接口跟踪修正
\begin{equation}
\min_{x\in\mathcal F}\frac{1}{H}\sum_{s,k}\pi_s\Delta t\left(\frac{|p_{s,k}-\hat p_{s,k}|}{P_{\rm base}}+\frac{|q_{s,k}-\hat q_{s,k}|}{Q_{\rm base}}\right)
\tag{20}
\end{equation}
\end{document}
